% Standalone section: exact positive-subcritical posterior lower bound.
\section{An exponential sample lower bound}
\label{sec:positive-sparse-partial}
Throughout this section $f=\chi_U$, where $U\subseteq[d]$ and $|U|=k$.
Write $D=d+1$, $\beta=2^{-(L-1)}$, and
\[
 B_{d,k}=\sqrt{k}\log\frac{eD}{\sqrt{k}},\qquad
 W_{d,k}=\max\{k,B_{d,k}\}.
\]
The $k$ term follows from the random-restriction amplitude tradeoff.
The second term is useful in very large ambient dimensions and follows
from a direct polynomial-threshold obstruction.

\begin{theorem}[Sparse parity for $0<\gamma<1$]
\label{thm:positive-sparse-partial}
Fix $L\ge2$ and $a>0$. There are constants $c>0$, $k_0$, and $d_0$,
depending only on $L,a$, such that the following holds. If
$d\ge d_0$, $k_0\le k\le d$, and a specified sample size satisfies
\[
 0\le n\le\exp(cW_{d,k}),
\]
then, with probability at least $3/4$ over independent uniform training
inputs, every minimizer of
\[
 \min\left\{\frac12\|z\|_K^2:
 |z(x)|\ge a\ \text{for all }x,\quad
 f(X_i)z(X_i)\ge a\ \text{for all }i\right\}
\]
has population sign error at least $3/8$.
On the same event, for every fixed $0<\gamma<1$, the actual network
posterior in the width-first limit satisfies
\[
 \liminf_{N\to\infty}
 \mathbb P_x\{q_{N,S}(x)\le2/3\}\ge\frac1{16}.
\]
The probability statement is for each specified sample size; it is not
a simultaneous event over all sample prefixes.
\end{theorem}

In particular, the lower bound is exponential in $k$ throughout the
entire growing sparse regime. It does not give the stronger scale
$\exp\{c k\log(ed/k)\}$ when $d/k\to\infty$.

\subsection{Squared balance on the target coordinates}
\begin{lemma}[Active-coordinate squared balance]
\label{lem:positive-active}
If $P$ has total Walsh degree at most $r$, then
\begin{equation}
 \mathbb E(fP)_-^2
 \ge\left(\frac12-\sqrt{\frac rk}\right)\|P\|_2^2.
 \label{eq:positive-active-balance}
\end{equation}
For every fixed $\delta>0$, there is a universal $A_\delta$ such that
if $r\le k/64$ and $n\ge A_\delta D_r\log(D_r+1)$, then, with
probability at least $1-\delta$,
\begin{equation}
 \mathbb P_n(fP)_-^2\ge\frac14\|P\|_2^2
 \quad\text{for every polynomial $P$ of degree at most $r$}.
 \label{eq:positive-active-sampling}
\end{equation}
\end{lemma}
\begin{proof}
For $j\in U$ put $D_jP=(P-P\circ\mathrm{flip}_j)/2$.
Because $f$ changes sign under these flips, the inequality
$|F(s)-F(t)|\le(|s|+|t|)|s-t|$ for $F(t)=t|t|$ gives
\[
 |\mathbb E fP|P||\le2\|P\|_2\|D_jP\|_2.
\]
Averaging over $j\in U$ and using
$\sum_{j\in U}\|D_jP\|_2^2\le r\|P\|_2^2$ proves
\eqref{eq:positive-active-balance} from
$(fP)_-^2=(P^2-fP|P|)/2$.
The vector of all Walsh characters of degree at most $r$ is isotropic
and has squared length $D_r$ at every input. Under $r\le k/64$,
the population bound is at least $3\|P\|_2^2/8$.
Apply Lemma~\ref{lem:squared-sampling} at accuracy $1/8$ and use
homogeneity to obtain \eqref{eq:positive-active-sampling}.
\end{proof}

\subsection{The polynomial-threshold range}
\begin{lemma}[A $\sqrt{k}$ cutoff]
\label{lem:positive-sqrt-range}
For fixed $L,a$, there are $c_b>0$, $k_b$, and $d_b$ such that, on the
localized-comparator event, every minimizer has error at least $3/8$
whenever $d\ge d_b$, $k_b\le k\le d$, and
\[
 n\le\exp\left(c_b\sqrt{k}\log\frac{eD}{\sqrt{k}}\right).
\]
\end{lemma}
\begin{proof}
The comparator gives the simultaneous bound
\begin{equation}
 \|z\|_K^2\le C_{L,a}(d+n d^\beta)
 \label{eq:positive-comparator-budget}
\end{equation}
for every minimizer. Choose a sufficiently small constant $\kappa>0$
and an odd $s$ with $s\sim\kappa\sqrt{k}$.
Fixing the $d-k$ nuisance coordinates leaves a degree-$s$ polynomial
on the active $k$-cube. Lemma~\ref{lem:ptf-parity} and the central
binomial bound imply that every such polynomial threshold has error
at least $7/16$, for all sufficiently large $k$.

The sharp spectrum gives
\[
 \|z-z_{\le s}\|_2^2
 \le C_{L,a}(d+n d^\beta)\lambda_{s+1}
 \le\frac{C_{L,a}(d+n d^\beta)s^{-1-\beta}}
            {\binom D{s+2}}.
\]
The binomial lower bound shows
\[
 \log\binom D{s+2}
 \ge c_1\kappa\sqrt{k}\log\frac{eD}{\sqrt{k}}
\]
for a universal $c_1>0$ once $k$ is large.
Choose $c_b<c_1\kappa/4$, then increase $k_b$ to absorb the factor
$d^\beta$ and increase $d_b$ to make the displayed tail bound at most
$a^2/16$, uniformly in the stated range. The additive $d$ is harmless
because $s$ is then a sufficiently large fixed integer or larger.
Since $|z|\ge a$, its sign differs from a fixed-tie threshold of
$z_{\le s}$ on at most $1/16$ of the cube. Its error is therefore
at least $7/16-1/16=3/8$.
The comparator's allowed sample range contains the present one for all
large $d$, because $\sqrt{k}\log(eD/\sqrt{k})=o(d)$ uniformly in $k\le d$.
\end{proof}

\subsection{Proof of the positive-subcritical theorem}
\begin{proof}[Proof of Theorem~\ref{thm:positive-sparse-partial}]
Fix $\delta=1/40$, choose $A$ large enough for
Lemma~\ref{lem:positive-active}, and fix an odd $r_0\ge11$.
Choose the small degree fraction in
Lemma~\ref{lem:sparse-gram} after these choices.
The localized comparator has failure probability $e^{-\Omega_L(d)}$.
All its conclusions hold simultaneously over the minimizing fields.

If $W_{d,k}=B_{d,k}$, choose $c\le c_b$ and apply
Lemma~\ref{lem:positive-sqrt-range}. We may therefore assume
$W_{d,k}=k$ and $n\le e^{ck}$.
Let $r$ be the largest odd integer satisfying
\[
 A D_r\log(D_r+1)\le n.
\]
If this cutoff does not exist or $r<r_0$, then
$n<A D_{r_0}\log(D_{r_0}+1)$.
After increasing $k_0$, this entire range lies below the threshold of
Lemma~\ref{lem:positive-sqrt-range}, which again proves the claim.

Suppose $r\ge r_0$. Choose $c<(\log2)/4$.
Since $k\le d$ and $D_r\ge2^{d-1}$ for $r\ge d/2$, the defining
lower bound on $n$ excludes $r\ge d/2$ for sufficiently large $d$.
In this range, the binomial lower bound implies
\begin{equation}
 rH\le C_0\log n\le C_0 c k,
 \qquad \rho=\frac rD,\quad H=\log\frac e\rho,
 \label{eq:positive-main-cutoff}
\end{equation}
for a universal constant $C_0$.
By taking $c$ sufficiently small, we ensure $r\le k/64$,
$r\le\alpha d$, and $n\le e^{d/25}$.
All comparator and sparse-Gram hypotheses therefore hold.
Moreover $r\ge r_0$ makes $n$ polynomially larger than $d$, so
\eqref{eq:positive-comparator-budget} becomes
\begin{equation}
 B:=\|z\|_K^2\le C_{L,a}n d^\beta.
 \label{eq:positive-main-budget}
\end{equation}

Work on the intersection of the comparator event, the squared-sampling
event \eqref{eq:positive-active-sampling}, and the Gram event in
Lemma~\ref{lem:sparse-gram}. Their total failure probability is
at most $2\delta+e^{-\Omega_L(d)}<1/4$ for sufficiently large $d$.
For an arbitrary minimizer write
\[
 z=P+b+t,\qquad P=z_{\le r},\quad b=z_{(r,r+2]},\quad t=z_{>r+2}.
\]
At a training point where $fP<0$, consistency gives
$(fP)_-\le|z-P|$. The sampling bound, RKHS evaluation bound, and
\eqref{eq:sparse-coarse-gram} yield
\[
 \frac14\|P\|_2^2
 \le\mathbb P_n(fP)_-^2
 \le\mathbb P_n(z-P)^2
 \le\frac Bn\|(Q_r)_{SS}\|_{\mathrm{op}}.
\]
Using \eqref{eq:positive-main-budget}, $\tau_r\asymp r^{-\beta}$,
and the deterministic estimates
$n\lambda_{r+1}\le C_L A H\tau_r$ and
$n\lambda_{r+3}\le C_L A\rho^2H\tau_r$, we obtain, simultaneously
for every minimizer,
\begin{align}
 \|P\|_2^2&\le C_{L,a,A,\delta}\rho^{-(\beta+1)}H,
 \label{eq:positive-partial-head}\\
 \|b\|_2^2&\le B\lambda_{r+1}
             \le C_{L,a,A}\rho^{-\beta}H,
 \label{eq:positive-partial-band}\\
 \|t\|_2^2&\le B\lambda_{r+3}
             \le C_{L,a,A}\rho^{2-\beta}H.
 \label{eq:positive-partial-tail}
\end{align}
For the first estimate we bounded
$1+\sqrt{A\rho^{-1}H}+AH$ by $C_A\rho^{-1}H$.
The eigenvalue estimates follow from the sharp spectrum and the
maximality interval for $r$; the second is also stated in
\eqref{eq:sparse-fine-eigenvalue}.
Put $h=P+b=z_{\le u}$, where $u=r+2\le2r$.
The first two pieces are orthogonal in $L^2$, and their sum therefore
has squared norm at most $C\rho^{-(\beta+1)}H$. Hence
Lemma~\ref{cor:restriction-amplitude} gives
\begin{equation}
 \mathbb P\{\operatorname{sign}z\ne f\}
 \ge\frac7{16}
       -C_1\rho^{2-\beta}H
       -C_1\rho^{-(\beta+1)}H\exp\!\left(-c_*\frac{k}{u}\right),
 \label{eq:positive-restriction-risk}
\end{equation}
where $c_*>0$ is universal and $C_1$ depends only on the already fixed
constants $L,a,A,\delta$.

We now choose the theorem's $c$. From
\eqref{eq:positive-main-cutoff}, $\rho\le C_0c$ and
$k/u\ge H/(2C_0c)$.
Since $x^{2-\beta}\log(e/x)\to0$ as $x\downarrow0$, the first error
term in \eqref{eq:positive-restriction-risk} is at most $1/32$ once
$c$ is sufficiently small. For the second term, use
$\rho^{-(\beta+1)}=\exp[(\beta+1)(H-1)]$ to obtain
\[
 C_1\rho^{-(\beta+1)}H e^{-c_*k/u}
 \le C_1 H\exp\!\left[
       \left(\beta+1-\frac{c_*}{2C_0c}\right)H\right].
\]
As $H\ge1$, taking $c$ still smaller makes this at most $1/32$
uniformly. These choices also ensure all earlier degree-fraction
restrictions. Finally increase $k_0,d_0$ to absorb finite rounding and
uniform asymptotic estimates. Thus every minimizing sign has error at
least $7/16-1/32-1/32=3/8$.

For completeness, let $q(x)$ be any probability mixture of these
minimizing signs, let $R=\mathbb E_x(1-q(x))$, and put
$F=\mathbb P_x\{q(x)\le2/3\}$.
The draw-error bound gives $R\ge3/8$, while
\[
 R\le\frac{1-F}{3}+F,
 \qquad F\ge\frac{3R-1}{2}\ge\frac1{16}.
\]
The positive-gamma width-limit proposition says that, for each fixed
$0<\gamma<1$, the actual finite-width posterior puts mass $1-o(1)$
on these minimum-energy signs. Thus its draw error is at least
$3/8-o(1)$, independently of the relative weights among ties.
Apply the same inequality at every width and take the lower limit.
This proves the asserted actual-posterior guarantee.
\end{proof}
